\clearpage
\VerdierPhysicalPageAnchor{0219}%
\begingroup\small\setlength{\parskip}{0pt}\setlength{\emergencystretch}{1.5em}
\setlength{\abovedisplayskip}{4pt plus 2pt minus 2pt}%
\setlength{\belowdisplayskip}{4pt plus 2pt minus 2pt}%

\VerdierSemanticAnchor{ch3d-3-3-2-sfl-i-induced-square}
où les lignes sont exactes fournit, pour tout entier $n$, un diagramme
commutatif :
\begin{center}
\begin{tikzcd}[column sep=5.2em,row sep=3.0em]
F^n(Z) \arrow[r,"{\delta_F^n(S)}"] \arrow[d,"{F^n(g)}"']
  & F^{n+1}(X) \arrow[d,"{F^{n+1}(f)}"] \\
F^n(Z') \arrow[r,"{\delta_F^n(S')}"']
  & F^{n+1}(X')
\end{tikzcd}
\end{center}

\VerdierSemanticAnchor{ch3d-3-3-2-sfl-ii}
\noindent\textbf{SFL II)} Lorsque
$S=0\longrightarrow0\longrightarrow0\longrightarrow0\longrightarrow0$,
le morphisme $\delta_F^n(S)$ est, pour tout entier $n$, un isomorphisme.

\VerdierSemanticAnchor{ch3d-3-3-2-msfl}
Soit maintenant $m\colon F\longrightarrow F'$ un morphisme de foncteurs
$F,F'\colon\mathcal A^{\operatorname{Ext}}\longrightarrow\mathcal C$.
On en déduit de façon évidente une suite de morphismes de foncteurs
$m^n\colon F^n\longrightarrow {F'}^n$. Cette suite de morphismes possède la
propriété :

\noindent\textbf{MSFL)} Pour toute suite exacte
$S=0\longrightarrow X\longrightarrow Y\longrightarrow Z\longrightarrow0$,
et pour tout entier $n$, le diagramme ci-après est commutatif :

\VerdierSemanticAnchor{ch3d-3-3-2-msfl-square}
\begin{center}
\begin{tikzcd}[column sep=5.2em,row sep=3.0em]
F^n(Z) \arrow[r,"{\delta_F^n(S)}"] \arrow[d,"{m^n(Z)}"']
  & F^{n+1}(X) \arrow[d,"{m^{n+1}(X)}"] \\
{F'}^n(Z) \arrow[r,"{\delta_{F'}^n(S)}"']
  & {F'}^{n+1}(X)
\end{tikzcd}
\end{center}

\VerdierSemanticAnchor{ch3d-3-3-3}
\medskip\noindent\textbf{Définition 3.3.3.}\quad
Une suite de foncteurs $F^n\colon\mathcal A\longrightarrow\mathcal C$ munie
de la donnée, pour toute suite exacte
$S=0\longrightarrow X\longrightarrow Y\longrightarrow Z\longrightarrow0$,
d'un morphisme
$\delta_F^n(S)\colon F^n(Z)\longrightarrow F^{n+1}(X)$ possédant les
propriétés SFL I) et SFL II) est appelée une \emph{suite de foncteurs liés}.
Soient $(F^n,\delta_F^n)$ et $({F'}^n,\delta_{F'}^n)$ deux suites de
foncteurs liés. Une suite de morphismes de foncteurs
$m^n\colon F^n\longrightarrow {F'}^n$ est appelée un \emph{morphisme de
suites de foncteurs liés} si elle possède la propriété MSFL).

\VerdierSemanticAnchor{ch3d-3-3-4}
\medskip\noindent\textbf{Proposition 3.3.4.}\quad
\textit{Soient
$\mathcal H\!om(\mathcal A^{\operatorname{Ext}},\mathcal C)$ la catégorie
des foncteurs de $\mathcal A^{\operatorname{Ext}}$ dans $\mathcal C$ et
$\mathcal{SFL}(\mathcal A,\mathcal C)$ la catégorie dont les objets sont les
suites de foncteurs liés et les morphismes, les morphismes de suites de
foncteurs liés. Ce qui précède définit un foncteur :}

\VerdierSemanticAnchor{ch3d-3-3-4-functor}
\[
\Phi\colon
\mathcal H\!om(\mathcal A^{\operatorname{Ext}},\mathcal C)
\longrightarrow\mathcal{SFL}(\mathcal A,\mathcal C).
\]

\VerdierSemanticAnchor{ch3d-3-3-4-a}
\noindent\textit{a) Ce foncteur est une équivalence de catégories.}
\endgroup
